\clearpage % Rozdziały zaczynamy od nowej strony.
\section{Moduł sztucznej inteligencji i analiza semantyczna obrazów}

Rozdział opisuje warstwę AI Workera, która rozpoznaje duplikaty
percepcyjne. Są to pliki o innej zawartości bajtowej, ale o tej samej
treści wizualnej, np. ponownie skompresowane, przeskalowane lub
przycięte kopie. Funkcje skrótu kryptograficznego nie wykrywają takich
par, ponieważ zmiana kodowania obrazu daje inny skrót
pliku~\cite{pizzi2022sscd}.

Moduł składa się z trzech części: ekstrakcji wektora cech modelem SSCD,
wyszukiwania najbliższego sąsiada w PostgreSQL z rozszerzeniem
\texttt{pgvector} oraz bezreferencyjnej oceny jakości. Podstawy
teoretyczne, czyli GeM, funkcję straty InfoNCE, regularyzację KoLeo
i miarę kosinusową, przedstawiono w rozdziale~2. Transport gRPC,
buforowanie strumienia i mapowanie statusów w bramie opisano
w podrozdziale~\ref{subsec:similarity-grpc-pipeline}. Poniżej
uwzględniono tylko elementy potrzebne do zrozumienia decyzji
podejmowanej przez AI Workera.

W całym rozdziale fakty o implementacji pochodzą z kodu repozytorium
\texttt{s3-dedup-engine}. Wnioski, które nie zostały zmierzone,
oznaczono jako przypuszczenia.

% ------------------------------------------------------------------
\subsection{Ekstrakcja cech i model SSCD (Self-Supervised Descriptor)}
\label{sec:model-sscd}

\subsubsection{Wykrywanie kopii a klasyfikacja obrazów}

Zadanie wykrywania kopii różni się od klasyfikacji. Autorzy SSCD
przyjmują, że dwa obrazy są kopiami wtedy i tylko wtedy, gdy pochodzą
z tego samego dwuwymiarowego źródła~\cite{pizzi2022sscd}.
Dwa różne zdjęcia tego samego psa nie są więc kopiami, mimo że należą
do tej samej klasy. Wektor odpowiedni dla klasyfikacji powinien
zbliżać obrazy jednej klasy. Wektor do wykrywania kopii powinien
natomiast utrzymywać w oddaleniu nawet obrazy semantycznie podobne,
a całą przestrzeń wypełniać możliwie równomiernie~\cite{pizzi2022sscd}.

Z tego powodu sieci wytrenowane z nadzorem na zbiorze
ImageNet~\cite{deng2009imagenet} nie są dobrym wyborem do tego zadania.
Uczą się one niezmienniczości względem cech, które nie zmieniają klasy
obiektu. W deduplikacji ta sama niezmienniczość prowadziłaby do
łączenia różnych zdjęć o podobnej treści. Tabela~\ref{tab:sscd-uap}
zestawia wyniki z pracy SSCD na zbiorze DISC2021.
Miara \(\mu\mathrm{AP}\) jest polem pod krzywą precyzja--czułość
wyznaczoną dla wszystkich par zwróconych przez system, łącznie
dla wszystkich zapytań~\cite{douze2021disc}. Odpowiada to stosowaniu
jednego, wspólnego progu odległości dla całego
zbioru~\cite{pizzi2022sscd}, czyli warunkom pracy bramy.

\begin{table}[!htbp]
\centering
\caption{Skuteczność wybranych deskryptorów w wykrywaniu kopii na
zbiorze DISC2021. \(\mu\mathrm{AP}_{SN}\) oznacza wynik po normalizacji
podobieństwa zbiorem tła. Opracowanie własne na podstawie
\cite{pizzi2022sscd}.}
\label{tab:sscd-uap}
\begin{tabular}{|l|r|r|r|}
\hline
\textbf{Model} & \textbf{Wymiar} & \(\boldsymbol{\mu\mathrm{AP}}\) [\%]
& \(\boldsymbol{\mu\mathrm{AP}_{SN}}\) [\%] \\
\hline
Multigrain (ImageNet, nadzorowany) & 1500 & 16,5 & 36,5 \\
Multigrain (ImageNet, nadzorowany) & 2048 & 20,5 & 41,7 \\
DINO ViT-B/16 & 1500 & 32,2 & 53,8 \\
SimCLR, wyjście sieci bazowej & 2048 & 13,1 & 33,9 \\
SimCLR, warstwa rzutująca & 128 & 9,4 & 17,3 \\
SSCD (DISC, adv.+mixup) & 512 & 61,5 & 72,5 \\
\hline
\end{tabular}
\end{table}

Model SSCD osiąga wyższy wynik od pozostałych deskryptorów przy
mniejszym wymiarze wektora niż modele Multigrain. Autorzy podają, że
SSCD przewyższa SimCLR o 48 punktów procentowych
\(\mu\mathrm{AP}\)~\cite{pizzi2022sscd}. Mniejszy wymiar ma znaczenie
praktyczne: zmniejsza rozmiar wiersza w bazie wektorowej i koszt
porównania pary wektorów.

\subsubsection{Budowa modelu}

SSCD wykorzystuje sieć ResNet-50~\cite{he2016resnet} jako sieć bazową.
Mapa cech z ostatniego bloku jest agregowana operacją GeM z
parametrem \(p=3\)~\cite{radenovic2018gem,pizzi2022sscd}. Następnie
stosowana jest warstwa liniowa rzutująca wektor do 512 wymiarów
i normalizacja L2. W przeciwieństwie do SimCLR~\cite{chen2020simclr}
warstwa rzutująca nie jest odrzucana po treningu, tylko pozostaje
częścią potoku wnioskowania~\cite{pizzi2022sscd}. Pełna ścieżka
wnioskowania ma więc postać:

\begin{equation}
\label{eq:sscd-forward}
\mathbf{e}(I)=
\frac{W\,\mathrm{GeM}_{p=3}\!\left(F(I)\right)}
{\left\|W\,\mathrm{GeM}_{p=3}\!\left(F(I)\right)\right\|_2},
\qquad \mathbf{e}(I)\in\mathbb{R}^{512},
\end{equation}

gdzie \(F\) oznacza sieć bazową ResNet-50, a \(W\) macierz warstwy
rzutującej. GeM uśrednia każdy kanał po wszystkich położeniach
przestrzennych, więc długość wektora nie zależy od wysokości i
szerokości mapy cech~\cite{radenovic2018gem}. Pozwala to podawać
na wejście obrazy o różnych proporcjach boków bez kadrowania.

Model był trenowany metodą samonadzorowaną, bez etykiet klas.
Pary pozytywne powstają jako dwie losowe transformacje tego samego
obrazu~\cite{chen2020simclr,pizzi2022sscd}. W wariancie
\texttt{sscd\_disc\_mixup} zestaw rozszerzony
(ang. \textit{advanced augmentations}) uzupełniono o mieszanie
obrazów metodami mixup i cutmix, stosowanymi z prawdopodobieństwem
po 2,5\%~\cite{pizzi2022sscd}. Wklejenie fragmentu jednego obrazu
w drugi modeluje kopie częściowe, np. kolaże~\cite{pizzi2022sscd}.
Zbiór DISC2021 zawiera właśnie takie przekształcenia: nakładki
tekstu i emotikon, zmiany kolorów, rozmycie, kompresję JPEG,
kadrowanie, obroty, odbicia oraz osadzanie obrazu w zrzucie ekranu
aplikacji~\cite{douze2021disc}. Odpowiadają one modyfikacjom,
którym podlegają obrazy przesyłane ponownie do magazynu obiektowego.

\subsubsection{Implementacja w AI Workerze}

Model ładowany jest z pliku TorchScript
\texttt{sscd\_disc\_mixup.torchscript.pt}, pobieranego z repozytorium
Hugging Face \texttt{m3/sscd-copy-detection}~\cite{hf_m3_sscd}.
Repozytorium to nie jest publikowane przez autorów SSCD; jego opis
odsyła do oficjalnego repozytorium Meta~\cite{hf_m3_sscd,sscd_repo},
w którym udostępniono model o tej samej nazwie: ResNet-50,
rozszerzenia adv.+mixup, wektor 512-wymiarowy~\cite{sscd_repo}.
Funkcja \texttt{hf\_hub\_download} jest wywoływana bez parametru
\texttt{revision}. Wersja pliku nie jest więc przypięta do konkretnej
rewizji repozytorium. Dla powtarzalności wyników należałoby zapisać
skrót pliku wag lub identyfikator rewizji.

Model wczytywany jest funkcją \texttt{torch.jit.load} na procesor CPU
i przełączany w tryb ewaluacji. Obiekt \texttt{SscdEmbedder} działa
jako singleton procesu, a pierwsze załadowanie chroni blokada
\texttt{threading.Lock}. Serwer wywołuje ładowanie przed przyjęciem
pierwszego żądania, dzięki czemu czas pobrania wag nie wlicza się
do obsługi pierwszego obrazu.

Przetwarzanie wstępne odpowiada zaleceniom autorów modelu.
Repozytorium SSCD zaleca skalowanie krótszego boku do 288 pikseli
albo skalowanie do kwadratu \(320\times320\)~\cite{sscd_repo}.
W pracy SSCD trening prowadzono w rozdzielczości \(224\times224\),
a ocenę po przeskalowaniu krótszego boku do 288 pikseli z zachowaniem
proporcji~\cite{pizzi2022sscd}. AI Worker stosuje pierwszy wariant.
Dla obrazu o wymiarach \(w\times h\) docelowy rozmiar wynosi:

\begin{equation}
\label{eq:sscd-resize}
s=\frac{288}{\min(w,h)},\qquad
(w',h')=\bigl(\max(1,\operatorname{round}(s\,w)),\;
\max(1,\operatorname{round}(s\,h))\bigr).
\end{equation}

Następnie każdy piksel kanału \(c\in\{R,G,B\}\) jest skalowany do
przedziału \([0,1]\) i standaryzowany statystykami ImageNet:

\begin{equation}
\label{eq:sscd-normalize}
\tilde{x}_c=\frac{x_c/255-\mu_c}{\sigma_c},\qquad
\boldsymbol{\mu}=(0{,}485;\,0{,}456;\,0{,}406),\quad
\boldsymbol{\sigma}=(0{,}229;\,0{,}224;\,0{,}225).
\end{equation}

Tabela~\ref{tab:sscd-preprocessing} zestawia kroki przetwarzania
wraz z miejscem ich realizacji w kodzie.

\begin{table}[!htbp]
\centering
\caption{Przetwarzanie obrazu przed wyznaczeniem wektora SSCD.
Opracowanie własne na podstawie implementacji.}
\label{tab:sscd-preprocessing}
\begin{tabular}{|p{0.27\textwidth}|p{0.63\textwidth}|}
\hline
\textbf{Krok} & \textbf{Realizacja w AI Workerze} \\
\hline
Dekodowanie & Pillow, konwersja do RGB; obsługiwane są formaty
rozpoznawane przez \texttt{Image.open}. \\
\hline
Skalowanie & Krótszy bok do 288 pikseli, interpolacja dwuliniowa,
bez kadrowania; równanie~\eqref{eq:sscd-resize}. \\
\hline
Normalizacja & Podział przez 255 i standaryzacja średnią i odchyleniem
ImageNet; równanie~\eqref{eq:sscd-normalize}. \\
\hline
Układ tensora & NCHW, partia o rozmiarze 1, typ \texttt{float32}. \\
\hline
Wnioskowanie & \texttt{torch.inference\_mode()}, CPU. \\
\hline
Kontrola wyjścia & Sprawdzenie liczby wymiarów (512) i ponowna
normalizacja L2 z dolnym ograniczeniem normy \(10^{-12}\). \\
\hline
\end{tabular}
\end{table}

Według repozytorium SSCD modele zwracają już wektory znormalizowane
w normie L2~\cite{sscd_repo}. Ponowna normalizacja w kodzie
\(\mathbf{e}\leftarrow\mathbf{e}/\max(\|\mathbf{e}\|_2,10^{-12})\)
nie zmienia więc poprawnego wyniku, a zabezpiecza przed zmianą pliku
wag na wariant bez normalizacji. Dolne ograniczenie normy chroni
przed dzieleniem przez zero.

Implementacja nie stosuje normalizacji podobieństwa zbiorem tła
opisanej w pracy SSCD. W tej metodzie od podobieństwa pary
odejmuje się podobieństwo zapytania do \(n\)-tego sąsiada
z oddzielnego zbioru obrazów niezwiązanych~\cite{pizzi2022sscd}:

\begin{equation}
\label{eq:sscd-score-norm}
\hat{s}(q,r)=s(q,r)-\beta\,s(q,b_n).
\end{equation}

Według tabeli~\ref{tab:sscd-uap} normalizacja podnosi
\(\mu\mathrm{AP}\) modelu SSCD z 61,5\% do 72,5\%. Jej wdrożenie
wymagałoby jednak utrzymywania i przeszukiwania drugiego indeksu
dla każdego obrazu. W bramie porównywana jest wyłącznie surowa
odległość kosinusowa.

% ------------------------------------------------------------------
\subsection{Reprezentacja wektorowa i przeszukiwanie przestrzeni
wielowymiarowej}
\label{sec:reprezentacja-wektorowa}

\subsubsection{Miara odległości}

Każdy obraz reprezentowany jest wektorem jednostkowym
\(\mathbf{e}\in\mathbb{R}^{512}\). Zależność między podobieństwem
kosinusowym a odległością zwracaną przez operator \texttt{<=>}
podano w równaniu~\eqref{eq:tech-cosine-distance}. Dla wektorów
jednostkowych zachodzi ponadto:

\begin{equation}
\label{eq:cos-euclid}
\|\mathbf{x}-\mathbf{y}\|_2^2
=\|\mathbf{x}\|_2^2+\|\mathbf{y}\|_2^2-2\,\mathbf{x}\cdot\mathbf{y}
=2-2\cos(\mathbf{x},\mathbf{y})
=2\,d_{\cos}(\mathbf{x},\mathbf{y}).
\end{equation}

Odległość kosinusowa jest więc połową kwadratu odległości
euklidesowej. Ranking sąsiadów według \(d_{\cos}\) jest taki sam jak
ranking według odległości euklidesowej. Pozwala to sprawdzić
spójność progów podawanych przez autorów modelu. Dla
\texttt{sscd\_disc\_mixup} repozytorium podaje równoważne progi:
podobieństwo kosinusowe powyżej 0,75, odległość euklidesowa poniżej
0,7 albo jej kwadrat poniżej 0,5~\cite{sscd_repo}.
Z równania~\eqref{eq:cos-euclid} wynika
\(0{,}7^2/2\approx0{,}245\) oraz \(0{,}5/2=0{,}25\), czyli
w przybliżeniu \(d_{\cos}<0{,}25\). Trzy postacie progu są zgodne.

Z równania~\eqref{eq:cos-euclid} wynika też, że \(d_{\cos}\)
nie jest metryką. Kwadrat odległości euklidesowej nie spełnia
nierówności trójkąta. Z warunków
\(d_{\cos}(\mathbf{a},\mathbf{b})\leq t\)
i \(d_{\cos}(\mathbf{b},\mathbf{c})\leq t\) nie wynika
\(d_{\cos}(\mathbf{a},\mathbf{c})\leq t\), a jedynie
\(d_{\cos}(\mathbf{a},\mathbf{c})\leq 4t\). Relacja „jest
duplikatem” zdefiniowana progiem nie jest zatem przechodnia.
Ma to znaczenie przy podmianie reprezentanta grupy, opisanej
w podrozdziale~\ref{sec:infogain-jakosc}.

\subsubsection{Przechowywanie wektorów}

AI Worker tworzy przy starcie rozszerzenie \texttt{vector}, tabelę
\texttt{image\_embeddings} i indeks HNSW. Schemat tabeli przedstawiono
w tabeli~\ref{tab:pg-schema}.

\begin{table}[!htbp]
\centering
\caption{Schemat tabeli \texttt{image\_embeddings}.
Opracowanie własne na podstawie implementacji.}
\label{tab:pg-schema}
\begin{tabular}{|l|l|p{0.45\textwidth}|}
\hline
\textbf{Kolumna} & \textbf{Typ} & \textbf{Znaczenie} \\
\hline
\texttt{id} & \texttt{BIGSERIAL}, klucz główny & Identyfikator wiersza,
używany przy blokadzie i podmianie. \\
\hline
\texttt{image\_key} & \texttt{TEXT}, \texttt{UNIQUE} & Nazwa pliku
reprezentanta grupy w magazynie. \\
\hline
\texttt{embedding} & \texttt{VECTOR(512)} & Wektor SSCD. \\
\hline
\texttt{quality\_score} & \texttt{DOUBLE PRECISION} & Wskaźnik jakości
z podrozdziału~\ref{sec:infogain-jakosc}. \\
\hline
\texttt{created\_at} & \texttt{TIMESTAMPTZ} & Czas utworzenia wiersza,
domyślnie \texttt{NOW()}. \\
\hline
\end{tabular}
\end{table}

Według dokumentacji \texttt{pgvector} wektor zajmuje
\(4d+8\) bajtów~\cite{pgvector_docs}. Dla \(d=512\) daje to
2056 bajtów na wektor, bez narzutu wiersza i indeksu.
Wektor przekazywany jest do zapytania jako tekst z ośmioma miejscami
po przecinku dla każdej składowej. Błąd zaokrąglenia składowej nie
przekracza \(5\cdot10^{-9}\). Dla wektora jednostkowego
\(\sum_i|x_i|\leq\sqrt{512}\), więc zmiana iloczynu skalarnego
nie przekracza około \(1{,}1\cdot10^{-7}\). Jest to wartość
pomijalna wobec progów decyzyjnych rzędu \(10^{-1}\).

W tabeli przechowywany jest tylko jeden wektor na grupę duplikatów.
Wektory obrazów odrzuconych jako duplikaty nie są zapisywane. Przy
podmianie wiersz reprezentanta jest aktualizowany w miejscu: zmienia
się nazwa, wektor i wskaźnik jakości. Liczba wierszy rośnie więc tylko przy statusie \texttt{inserted}
z nową nazwą pliku, a nie zmienia się przy odrzuceniach i podmianach.

\subsubsection{Indeks HNSW}

Indeks utworzono metodą \texttt{hnsw} z klasą operatorów
\texttt{vector\_\allowbreak cosine\_\allowbreak ops}, bez podawania
parametrów. Obowiązują więc wartości domyślne \texttt{pgvector}:
\(m=16\), \(\mathit{ef\_construction}=64\)
i \(\mathit{hnsw.ef\_search}=40\)~\cite{pgvector_docs}.

HNSW (ang. \textit{Hierarchical Navigable Small World}) buduje
wielowarstwowy graf bliskości. Wyższe warstwy zawierają coraz mniej
punktów i długie połączenia, a warstwa najniższa wszystkie punkty.
Wyszukiwanie zaczyna się od warstwy najwyższej i schodzi zachłannie
w dół, zawężając obszar poszukiwań~\cite{malkov2020hnsw}. Parametr
\(m\) ogranicza liczbę połączeń wierzchołka, \(\mathit{ef\_construction}\)
rozmiar listy kandydatów przy budowie grafu, a \(\mathit{ef\_search}\)
rozmiar tej listy przy zapytaniu~\cite{malkov2020hnsw,pgvector_docs}.
Według dokumentacji \texttt{pgvector} HNSW daje lepszy kompromis
między szybkością a czułością niż IVFFlat, kosztem dłuższej budowy
i większego zużycia pamięci~\cite{pgvector_docs}. Indeks można
też utworzyć na pustej tabeli, co odpowiada sposobowi inicjalizacji
schematu w AI Workerze.

Wyszukiwanie najbliższego sąsiada ma postać:

\begin{lstlisting}[basicstyle=\ttfamily\small,breaklines=true]
SELECT id, image_key, quality_score,
       embedding <=> %(vector)s::vector AS distance
FROM image_embeddings
ORDER BY distance
LIMIT 1;
\end{lstlisting}

Zapytanie zawiera \texttt{ORDER BY} odległości rosnąco i
\texttt{LIMIT}, czyli warunki, przy których planista może użyć indeksu
HNSW~\cite{pgvector_docs}. Dla małej tabeli planista może wybrać
skanowanie sekwencyjne; wynik jest wtedy dokładny. Przy użyciu indeksu
wyszukiwanie jest przybliżone: dokumentacja zaznacza, że indeks
zamienia część czułości na szybkość i może zwrócić inne wyniki niż
skanowanie pełne~\cite{pgvector_docs}.

Skutki przybliżenia dla decyzji można określić dokładnie. Niech
\(\mathbf{v}^{*}\) będzie rzeczywistym najbliższym sąsiadem, a
\(\mathbf{v}'\) sąsiadem zwróconym przez indeks. Odległość
\(d'=d_{\cos}(\mathbf{e},\mathbf{v}')\) jest obliczana dokładnie
dla zwróconego wiersza i spełnia \(d'\geq d^{*}\). Możliwe są dwa
przypadki błędne:

\begin{itemize}
    \item \(d^{*}\leq t<d'\): obraz zostaje przyjęty jako nowy, mimo że
    w bazie istnieje duplikat. Jest to błąd pominięcia.
    \item \(d^{*}\leq d'\leq t\), \(\mathbf{v}'\neq\mathbf{v}^{*}\):
    obraz porównywany jest z reprezentantem innej grupy, która również
    mieści się w progu.
\end{itemize}

Przybliżenie nie może natomiast spowodować uznania za duplikat
obrazu, którego zwrócony sąsiad leży dalej niż próg, ponieważ warunek
progowy sprawdzany jest na dokładnie obliczonej odległości.
Liczba zwracanych wierszy wynosi jeden, więc filtrowanie po stronie
zapytania i iteracyjne skanowanie indeksu, dostępne od wersji 0.8.0
\texttt{pgvector}~\cite{pgvector_docs}, nie mają tu zastosowania.

% ------------------------------------------------------------------
\subsection{Obiektywna ocena jakości obrazu (InfoGain) i logika
decyzyjna}
\label{sec:infogain-jakosc}

\subsubsection{Problem wyboru lepszej wersji}

Po znalezieniu sąsiada w progu odległości trzeba zdecydować, która
wersja pozostanie w magazynie. Nie można przyjąć, że plik przesłany
wcześniej jest oryginałem. Kolejna wersja może mieć wyższą
rozdzielczość albo mniejszą kompresję. Wskaźnik jakości musi przy tym
działać bez obrazu wzorcowego. Miary pełnoreferencyjne, takie jak
SSIM, porównują obraz z wersją uznaną za
niezniekształconą~\cite{wang2004ssim}. W deduplikacji żadna z wersji
nie jest takim wzorcem, a o jego wyborze ma dopiero rozstrzygnąć
porównanie.

Potrzebna jest więc miara bezreferencyjna. W literaturze stosuje się
do tego m.in. modele statystyk scen naturalnych, takie jak BRISQUE,
uczone na ocenach jakości wystawionych przez
ludzi~\cite{mittal2012brisque}. W projekcie wybrano prostszy
wskaźnik heurystyczny złożony z trzech cech, które można policzyć
bez dodatkowego modelu. Nazwa \textit{InfoGain} pochodzi z opisu
modułu \texttt{quality/score.py}, który określa wynik jako zysk
informacyjny deduplikacji. Nie jest to miara teorioinformacyjna,
lecz ważona suma cech obrazu.

\subsubsection{Definicja wskaźnika}

Obraz dekodowany jest raz. Najpierw odczytywane są wymiary
oryginalne \(w\times h\). Następnie tworzony jest obraz w skali
szarości, przeskalowany tak, aby dłuższy bok miał 512 pikseli.
Dla plików JPEG i MPO wywoływana jest metoda \texttt{Image.draft},
która konfiguruje dekoder tak, aby od razu zwrócił obraz możliwie
bliski żądanemu trybowi i rozmiarowi~\cite{pillow_image}. Ogranicza
to koszt dekodowania dużych zdjęć. Dokumentacja podaje, że metoda jest obsługiwana tylko dla tych
dwóch formatów~\cite{pillow_image}. Końcowe skalowanie do 512 pikseli
wykonywane jest interpolacją dwuliniową, także w górę dla obrazów
mniejszych.

Na tej podstawie liczone są trzy składniki. Każdy jest normalizowany
funkcją:

\begin{equation}
\label{eq:clamped-ratio}
\rho(v,r)=
\begin{cases}
0, & r\leq0,\\
\min\!\left(\max\!\left(\dfrac{v}{r},0\right),1\right), & r>0.
\end{cases}
\end{equation}

\paragraph{Rozdzielczość.}
Składnik ten porównuje liczbę pikseli oryginału, sprzed
zmniejszenia, z liczbą pikseli obrazu UHD:

\begin{equation}
\label{eq:q-resolution}
q_{\mathrm{res}}=\rho\!\left(w\,h,\;3840\cdot2160\right).
\end{equation}

Obrazy o liczbie pikseli równej lub większej od
\(8\,294\,400\) otrzymują wartość 1. Powyżej tej granicy
rozdzielczość nie różnicuje już wersji.

\paragraph{Ostrość.}
Ostrość mierzona jest wariancją laplasjanu obrazu w skali szarości
\(G\) o dłuższym boku 512 pikseli:

\begin{equation}
\label{eq:q-sharpness}
q_{\mathrm{sharp}}=\rho\!\left(\operatorname{Var}\!\left(\nabla^2 G\right),\;600\right).
\end{equation}

OpenCV wyznacza laplasjan jako sumę drugich pochodnych
cząstkowych. Przy domyślnym rozmiarze jądra stosowana jest maska
\(\begin{bmatrix}0&1&0\\1&-4&1\\0&1&0\end{bmatrix}\)~\cite{opencv_laplacian}.
Wariancję laplasjanu zaproponowano jako miarę ostrości do
automatycznego ustawiania ostrości mikroskopu~\cite{pechpacheco2000}.
W porównaniu operatorów miary ostrości operatory oparte na
laplasjanie dały najlepsze wyniki w typowych warunkach, ale były
najbardziej wrażliwe na szum~\cite{pertuz2013focus}. Obraz
zaszumiony może więc uzyskać zawyżoną ocenę ostrości. Rozmycie
obniża wariancję, ponieważ tłumi szybkie zmiany jasności.

Pomiar odbywa się na obrazie o stałym rozmiarze, więc ostrość jest
porównywalna między plikami o różnej rozdzielczości. Obraz mały,
powiększony do 512 pikseli przez interpolację, nie zyskuje nowych
szczegółów. Można się spodziewać, że jego wariancja laplasjanu będzie
niższa niż dla tej samej sceny zmniejszonej z wysokiej rozdzielczości.
Składnik ostrości jest wtedy częściowo skorelowany z rozdzielczością.
Tego przypuszczenia nie zweryfikowano pomiarowo.

\paragraph{Kontrast.}
Kontrast mierzony jest odchyleniem standardowym jasności pikseli
obrazu \(G\), zapisanych jako liczby całkowite z przedziału
\([0,255]\):

\begin{equation}
\label{eq:q-contrast}
q_{\mathrm{con}}=\rho\!\left(\sigma(G),\;80\right).
\end{equation}

Jest to globalny kontrast RMS. Peli wskazuje, że pojedyncza wartość
dla całego obrazu nie oddaje kontrastu postrzeganego lokalnie
w obrazach złożonych~\cite{peli1990contrast}. W tym zastosowaniu
składnik ma ograniczoną rolę: odróżnia wersje wyblakłe od wersji
o pełnym zakresie tonalnym.

\paragraph{Wynik łączny.}
Wskaźnik jest ważoną sumą składników:

\begin{equation}
\label{eq:quality-total}
Q=0{,}5\,q_{\mathrm{res}}+0{,}3\,q_{\mathrm{sharp}}+0{,}2\,q_{\mathrm{con}},
\qquad Q\in[0,1].
\end{equation}

Tabela~\ref{tab:quality-factors} zestawia parametry wskaźnika.

\begin{table}[!htbp]
\centering
\caption{Składniki wskaźnika jakości.
Opracowanie własne na podstawie implementacji.}
\label{tab:quality-factors}
\begin{tabular}{|l|p{0.33\textwidth}|r|r|}
\hline
\textbf{Składnik} & \textbf{Wielkość mierzona} &
\textbf{Odniesienie} & \textbf{Waga} \\
\hline
\(q_{\mathrm{res}}\) & liczba pikseli oryginału &
\(3840\cdot2160\) & 0,5 \\
\hline
\(q_{\mathrm{sharp}}\) & wariancja laplasjanu na obrazie 512 px &
600 & 0,3 \\
\hline
\(q_{\mathrm{con}}\) & odchylenie standardowe jasności &
80 & 0,2 \\
\hline
\end{tabular}
\end{table}

Wagi i wartości odniesienia są stałymi projektowymi w pliku
\texttt{const/quality.py}. W repozytorium nie ma procedury, która
dopasowywałaby je do ocen ludzi lub do zbioru testowego. Testy
jednostkowe w \texttt{tests/test\_quality\_score.py} sprawdzają
własności wskaźnika, a nie jego zgodność z percepcją:

\begin{itemize}
    \item wagi sumują się do jedności, a waga rozdzielczości jest
    największa;
    \item wynik maleje monotonicznie dla tej samej sceny w kolejności
    UHD, HD, SD i miniatura;
    \item przy tej samej rozdzielczości obraz ostrzejszy otrzymuje
    wyższy wynik niż rozmyty;
    \item każdy składnik leży w przedziale \([0,1]\), a wynik jest
    deterministyczny.
\end{itemize}

Wskaźnik nie zawiera osobnej miary artefaktów kompresji, np.
blokowości JPEG. Artefakty wpływają na wynik tylko pośrednio, przez
zmianę wariancji laplasjanu i rozkładu jasności. Krawędzie bloków
\(8\times8\) dodają składowe o wysokiej częstotliwości, więc mogą
nawet podnieść ocenę ostrości. Jest to wniosek z definicji
operatora, niezmierzony w projekcie.

\subsubsection{Przykład działania wag}

Waga rozdzielczości decyduje o tym, kiedy mniejsza wersja może wygrać
z większą. Nowy obraz \(n\) zastępuje istniejący obraz \(o\), gdy
\(Q_n>Q_o\). Po podstawieniu równania~\eqref{eq:quality-total}:

\begin{equation}
\label{eq:quality-tradeoff}
0{,}3\,\Delta q_{\mathrm{sharp}}+0{,}2\,\Delta q_{\mathrm{con}}
>0{,}5\,\left(q^{o}_{\mathrm{res}}-q^{n}_{\mathrm{res}}\right),
\end{equation}

gdzie \(\Delta q=q^{n}-q^{o}\). Przykładowo zdjęcie
\(4000\times3000\) ma \(q_{\mathrm{res}}=1\), a jego kopia
\(1280\times960\) ma \(q_{\mathrm{res}}\approx0{,}148\). Prawa strona
nierówności wynosi wtedy około 0,426. Lewa strona nie przekracza
0,5, więc mniejsza kopia wygra tylko wtedy, gdy będzie znacznie
ostrzejsza i bardziej kontrastowa od dużego, np. silnie rozmytego
oryginału. Dla dwóch obrazów co najmniej w rozdzielczości UHD prawa
strona wynosi zero i o wyniku decydują wyłącznie ostrość i kontrast.

\subsubsection{Logika decyzyjna}

Kolejność operacji w AI Workerze jest stała. Po odebraniu całego
strumienia obliczany jest wskaźnik \(Q\), następnie wektor
\(\mathbf{e}\), a na końcu wywoływana jest metoda
\texttt{Embedding\allowbreak Repository.\allowbreak resolve}.
Czasy etapów trafiają do dziennika w polach
\texttt{score\_ms}, \texttt{embed\_ms} i \texttt{resolve\_ms}. Decyzję można zapisać jako:

\begin{equation}
\label{eq:ai-decision}
\mathrm{status}=
\begin{cases}
\texttt{inserted}, & \text{brak sąsiada lub } d'>t,\\
\texttt{duplicate\_rejected}, & d'\leq t \;\wedge\; Q_n\leq Q_o,\\
\texttt{replaced}, & d'\leq t \;\wedge\; Q_n> Q_o,
\end{cases}
\end{equation}

gdzie \(t\) jest progiem przekazanym w polu \texttt{threshold},
domyślnie \(0{,}60\). Przy pustej tabeli zwracana jest odległość
umowna 1,0. Remis wskaźnika jakości prowadzi do odrzucenia nowego
obrazu, więc przy równych ocenach zachowywana jest wersja wcześniejsza.
Przebieg metody przedstawiono na rysunku~\ref{fig:ai-resolve-algo}.

\begin{figure}[!htbp]
\centering
\begin{lstlisting}[basicstyle=\ttfamily\small,breaklines=true]
BEGIN
  pg_advisory_xact_lock(74821001)
  nb = nearest(e)                    -- ORDER BY <=> LIMIT 1
  if nb is None or nb.distance > t:
      INSERT ... ON CONFLICT (image_key) DO UPDATE
      return inserted
  q_old = SELECT quality_score ... WHERE id = nb.id FOR UPDATE
  if q_new <= q_old:
      return duplicate_rejected(nb.image_key)
  UPDATE image_embeddings
     SET image_key, embedding, quality_score
   WHERE id = nb.id
  return replaced(nb.image_key)
COMMIT
\end{lstlisting}
\caption{Rozstrzygnięcie decyzji w \texttt{EmbeddingRepository.resolve}.
Opracowanie własne na podstawie implementacji.}
\label{fig:ai-resolve-algo}
\end{figure}

Cała metoda działa w jednej transakcji. Transakcyjna blokada doradcza
o stałym identyfikatorze szereguje rozstrzygnięcia wszystkich obrazów,
a \texttt{SELECT ... FOR UPDATE} blokuje wiersz reprezentanta do
końca transakcji~\cite{postgres16_locking}. Zakres tych gwarancji
i brak wspólnej transakcji z zapisami S3 i DynamoDB omówiono
w podrozdziale~\ref{subsec:wspolbieznosc}. Ocena jakości i wyznaczenie
wektora odbywają się przed pobraniem blokady, więc mogą przebiegać
równolegle w czterech wątkach serwera gRPC. Szeregowane jest tylko
wyszukanie i zapis.

Status przekazywany jest do bramy. W pliku \texttt{decision.go}
zdefiniowano trzy decyzje potoku: \texttt{continue},
\texttt{accept\_full} i \texttt{soft\_dedup}.
W pliku \texttt{similarity.go} funkcja
\texttt{apply\allowbreak Semantic\allowbreak Status} mapuje na nie statusy AI Workera zgodnie
z tabelą~\ref{tab:ai-status-mapping}.

\begin{table}[!htbp]
\centering
\caption{Mapowanie statusów AI Workera na działania bramy.
Opracowanie własne na podstawie implementacji.}
\label{tab:ai-status-mapping}
\begin{tabular}{|p{0.25\textwidth}|p{0.25\textwidth}|p{0.38\textwidth}|}
\hline
\textbf{Status} & \textbf{Decyzja bramy} & \textbf{Skutek w magazynie} \\
\hline
\texttt{inserted} & \texttt{accept\_full} & Zapis obrazu, JSON
i wpisów DynamoDB. \\
\hline
\texttt{duplicate\_\allowbreak rejected} & \texttt{soft\_dedup} & Pominięcie
zapisu obrazu; zapisywany jest JSON ze wskazaniem
\texttt{existing\_image\_key}. \\
\hline
\texttt{replaced} & \texttt{accept\_full} oraz flaga \texttt{Replace\allowbreak Existing}
& Zapis nowego obrazu, a następnie usunięcie poprzedniego obiektu,
jego JSON i wpisów DynamoDB. \\
\hline
\end{tabular}
\end{table}

Odrzucenie (\texttt{duplicate\_rejected}) i podmiana
(\texttt{replaced}) dają różne skutki kosztowe. Odrzucenie zmniejsza
liczbę zapisów do S3. Podmiana zwiększa liczbę operacji: obejmuje
zapis nowego obrazu, odczyt rozmiaru poprzedniego obiektu i jego
usunięcie. Względem stanu sprzed operacji objętość danych zmienia się
o różnicę rozmiarów obu plików i przy wersji lepszej jakości może
wzrosnąć. Względem systemu bez deduplikacji oszczędzany jest rozmiar
usuniętego obiektu. Bilans tych operacji przedstawiono
w rozdziale~\ref{chap:ewaluacja-i-finops}.

\subsubsection{Konsekwencje przyjętego modelu grupy}

Grupa duplikatów reprezentowana jest jednym wektorem. Po podmianie
wiersz otrzymuje wektor nowego obrazu, a wektor poprzedniego jest
tracony. Kolejne obrazy porównywane są więc z nowym reprezentantem.
Ponieważ relacja progowa nie jest przechodnia, obraz bliski dawnemu
reprezentantowi może leżeć poza progiem od nowego i zostać przyjęty
jako odrębny. Ten sam zbiór plików może zatem dać inny wynik
deduplikacji przy innej kolejności przesyłania.

Decyzja uwzględnia tylko jednego sąsiada. Jeżeli w progu mieści się
kilka reprezentantów, porównanie jakości odbywa się wyłącznie
z najbliższym z nich. Pozostałe grupy nie są łączone.

Operacja \texttt{INSERT ... ON CONFLICT (image\_key) DO UPDATE}
nadpisuje wiersz, gdy obraz o tej samej nazwie nie ma sąsiada
w progu. Ponowne przesłanie innego obrazu pod tą samą nazwą zastępuje
więc wektor i wskaźnik jakości poprzedniego pliku bez porównania
jakości.

% ------------------------------------------------------------------
\subsection{Ewaluacja i strojenie progów podobieństwa (Threshold Tuning)}
\label{sec:ewaluacja-progów}

\subsubsection{Miary oceny progu}

Próg \(t\) dzieli pary obrazów na uznane i nieuznane za duplikaty.
Niech \(P\) będzie zbiorem par rzeczywiście zduplikowanych, a \(N\)
zbiorem par różnych. Dla progu \(t\) definiuje się odsetek wykrytych
duplikatów i odsetek fałszywych alarmów~\cite{fawcett2006roc}:

\begin{equation}
\label{eq:tpr-fpr}
\mathrm{TPR}(t)=\frac{\left|\{(a,b)\in P: d_{\cos}(a,b)\leq t\}\right|}{|P|},
\qquad
\mathrm{FPR}(t)=\frac{\left|\{(a,b)\in N: d_{\cos}(a,b)\leq t\}\right|}{|N|}.
\end{equation}

Przebieg obu wielkości przy zmianie \(t\) tworzy krzywą
ROC~\cite{fawcett2006roc}. W deduplikacji koszty błędów są
niesymetryczne. Fałszywy alarm prowadzi do pominięcia zapisu albo
usunięcia innego obrazu, czyli do utraty danych. Pominięcie duplikatu
powoduje jedynie dodatkowy koszt przechowywania. Próg powinien więc
w pierwszej kolejności ograniczać FPR.

W repozytorium przygotowano dwa notatniki: \texttt{tune\_threshold.ipynb}
do analizy rozkładu odległości par oraz \texttt{eval\_meta.ipynb}
do oceny całego potoku na zbiorze DISC2021. Skrypt
\texttt{scripts/fetch\_dataset.py}, przewidziany do pobierania danych,
jest pusty. Zbiory trzeba więc przygotować ręcznie.

\subsubsection{Rozkład odległości na zbiorze bliskich duplikatów}

Notatnik \texttt{tune\_threshold.ipynb} korzysta ze zbioru
\textit{Near Duplicate Dataset} z serwisu Kaggle. Zbiór zawiera zdjęcia
produktów sklepowych wykonane pod nieco innym kątem lub
w innym czasie; jest udostępniony na licencji CC BY-NC
4.0~\cite{kaggle_crystal0801}. Opis zbioru podaje 455 plików JPG;
notatnik wczytał 456 plików w 64 grupach obiektów, rozpoznawanych
po nazwie \texttt{\{grupa\} (\{nr\}).jpg}. Wzorzec nazw w notatniku dopuszcza także
przyrostek \texttt{clone}. Przyczyny różnicy jednego pliku
nie sprawdzono.

Procedura notatnika jest następująca:

\begin{enumerate}
    \item Utworzenie wszystkich par w obrębie grupy (2113 par
    duplikatów) i między grupami (101\,627 par różnych).
    \item Losowanie po 200 par z każdego zbioru (ziarno 7).
    \item Wyznaczenie wektorów obiektem zwracanym przez
    \texttt{get\_embedder()} z modułu AI Workera, czyli tym samym
    modelem i przetwarzaniem wstępnym; łącznie dla 377 różnych zdjęć z wylosowanych par.
    \item Obliczenie odległości funkcją \texttt{cosine} z modułu
    \texttt{scipy.\allowbreak spatial.\allowbreak distance}. Funkcja zwraca
    \(1-\cos(\mathbf{x},\mathbf{y})\)~\cite{scipy_cosine}, czyli tę samą
    wielkość co operator \texttt{<=>}~\cite{pgvector_docs}.
\end{enumerate}

Wyniki przedstawiono w tabeli~\ref{tab:tune-threshold}.

\begin{table}[!htbp]
\centering
\caption{Rozkład odległości kosinusowej SSCD dla losowych par ze zbioru
\textit{Near Duplicate Dataset} (po 200 par).
Opracowanie własne na podstawie wyników notatnika
\texttt{tune\_threshold.ipynb}.}
\label{tab:tune-threshold}
\begin{tabular}{|l|r|r|r|r|}
\hline
\textbf{Rodzaj pary} & \textbf{Średnia} & \textbf{Mediana} &
\textbf{Percentyl 90} & \(\boldsymbol{d_{\cos}\leq0{,}10}\) \\
\hline
Ta sama grupa obiektu & 0,3455 & 0,3355 & 0,4832 & 1,5\% \\
Różne grupy & 0,8850 & 0,8939 & 0,9867 & 0\% \\
\hline
\end{tabular}
\end{table}

Przykładowe pary z jednej grupy miały odległości 0,2221, 0,3885
i 0,4438, a pary z różnych grup 0,7353 i 0,7714. Histogram
rozkładów nie został zapisany w notatniku.

Wyniki należy interpretować w świetle definicji kopii przyjętej
w SSCD. Zdjęcia w tym zbiorze przedstawiają ten sam obiekt, ale są
odrębnymi ujęciami, a nie kopiami jednego pliku. Według definicji
z pracy SSCD nie są to więc kopie~\cite{pizzi2022sscd}. Model nie był
trenowany do ich zbliżania. Wyjaśnia to, dlaczego mediana odległości
w grupie (0,3355) przekracza próg około 0,25 podany przez autorów
modelu dla precyzji 90\% na DISC2021~\cite{sscd_repo}. Przy tamtym
progu co najmniej połowa par z tabeli~\ref{tab:tune-threshold} nie
zostałaby uznana za duplikaty.

Z tabeli można wyprowadzić tylko część wniosków o progu domyślnym
\(t=0{,}60\). Percentyl 90 odległości par z jednej grupy wynosi
0,4832, więc co najmniej 90\% wylosowanych par spełnia
\(d_{\cos}\leq0{,}60\), czyli \(\mathrm{TPR}(0{,}60)\geq0{,}9\)
na tej próbie. Notatnik nie podaje natomiast odsetka par różnych
poniżej 0,60. Mediana 0,8939 ogranicza go jedynie od góry przez 50\%,
co nie wystarcza do oceny FPR. Odległości nie są zapisywane do pliku, więc brakującą wielkość
trzeba obliczyć po ponownym uruchomieniu notatnika.

Dla progu 0,10 notatnik wykazał 0 z 200 par różnych. Zero
w próbie nie oznacza zerowego ryzyka. Według reguły trzech górna
granica 95-procentowego przedziału ufności dla odsetka zdarzeń przy
zerze obserwacji w \(n\) próbach wynosi około
\(3/n\)~\cite{hanley1983rule}. Dla \(n=200\) daje to 1,5\%. Przy tym
samym progu wykryto jednak tylko 3 z 200 par duplikatów (1,5\%).

W opisie notatnika wspomniano stałą \texttt{MAX\_COSINE\_DISTANCE = 0.10}
w pliku \texttt{const/pg.py}. W bieżącym kodzie taka stała nie
istnieje. Próg przekazywany jest z bramy w polu \texttt{threshold}
i domyślnie wynosi 0,60 (\texttt{DefaultCosineThreshold}
w \texttt{internal/constants/const.go}). Opis notatnika jest więc
niezgodny z implementacją i dotyczy wcześniejszej wersji systemu.
Wartości 0,60 nie wyznaczono procedurą optymalizacji opisaną
w repozytorium. Jest to wartość domyślna, którą należy potwierdzić
pomiarem FPR na danych docelowych.

\subsubsection{Ocena potoku na zbiorze DISC2021}

Notatnik \texttt{eval\_meta.ipynb} ocenia cały system, a nie tylko
rozkład odległości. Wykorzystuje zbiór DISC2021, przygotowany przez
Meta do konkursu Image Similarity Challenge 2021~\cite{douze2021disc,
papakipos2022isc}. Zbiór zawiera milion obrazów referencyjnych bez
przekształceń oraz po 50\,000 zapytań w zbiorach rozwojowym
i testowym. W każdym z nich 10\,000 zapytań to przekształcone kopie
obrazów referencyjnych, a 40\,000 to obrazy rozpraszające bez
odpowiednika w zbiorze referencyjnym~\cite{douze2021disc}. Dla zbioru
rozwojowego udostępniono przyporządkowania dla 25\,000
zapytań~\cite{douze2021disc}.

Notatnik korzysta z jednego archiwum obrazów referencyjnych
(\texttt{refs\_50k\_0}), jednego archiwum zapytań rozwojowych
(\texttt{dev\_queries\_50k\_0}) oraz pliku
\texttt{dev\_queries\_groundtruth.csv}. Parametry przebiegu podano
w tabeli~\ref{tab:eval-meta-params}.

\begin{table}[!htbp]
\centering
\caption{Parametry notatnika \texttt{eval\_meta.ipynb}.
Opracowanie własne na podstawie implementacji.}
\label{tab:eval-meta-params}
\begin{tabular}{|l|r|p{0.5\textwidth}|}
\hline
\textbf{Parametr} & \textbf{Wartość} & \textbf{Znaczenie} \\
\hline
\texttt{THRESHOLD} & 0,60 & Próg przekazywany bramie
(\texttt{-{}-threshold}). \\
\hline
\texttt{WORKERS} & 10 & Liczba równoległych zadań bramy
(\texttt{-{}-workers}). \\
\hline
\texttt{LIMIT} & 200 & Liczba obrazów referencyjnych i liczba zapytań. \\
\hline
\texttt{DUP\_FRACTION} & 0,20 & Udział dokładnych kopii obrazów
referencyjnych (40 plików). \\
\hline
\texttt{BATCH\_SIZE} & 40 & Liczba plików w jednym uruchomieniu bramy. \\
\hline
\texttt{SEED} & 7 & Ziarno losowania i mieszania. \\
\hline
\end{tabular}
\end{table}

Procedura obejmuje następujące kroki:

\begin{enumerate}
    \item Wylosowanie 200 obrazów referencyjnych i 200 zapytań,
    dla których istnieje przyporządkowanie w pliku odpowiedzi.
    \item Utworzenie 40 dokładnych kopii pierwszych obrazów
    referencyjnych z przyrostkiem \texttt{\_dup} w nazwie.
    \item Wymieszanie całej listy 440 plików i podział na partie
    po 40 plików.
    \item Rozpakowanie partii i uruchomienie bramy poleceniem
    \texttt{go run ./cmd/gateway}. Przekazywana jest lista plików
    (\texttt{-{}-list}), próg 0,60 (\texttt{-{}-threshold})
    i liczba zadań 10 (\texttt{-{}-workers}).
    \item Odczyt zdarzeń z pliku \texttt{events.jsonl}.
    Każdy plik otrzymuje grupę (\texttt{ref}, \texttt{exact\_dup}
    lub \texttt{augmented}) oraz wynik bramy.
\end{enumerate}

Wynik bramy zapisany w polu \texttt{outcome} przyjmuje wartości
\texttt{inserted}, \texttt{hash\_\allowbreak duplicate},
\texttt{semantic\_\allowbreak duplicate},
\texttt{replaced} i \texttt{failed}. Dla grup \texttt{exact\_dup}
i \texttt{augmented} notatnik porównuje nazwę wskazaną w polu
\texttt{existing\_image\_key} z oczekiwanym obrazem referencyjnym.
Etykiety oceny opisano w tabeli~\ref{tab:eval-meta-labels}.

\begin{table}[!htbp]
\centering
\caption{Etykiety oceny w notatniku \texttt{eval\_meta.ipynb}.
Opracowanie własne na podstawie implementacji.}
\label{tab:eval-meta-labels}
\begin{tabular}{|l|p{0.62\textwidth}|}
\hline
\textbf{Etykieta} & \textbf{Warunek} \\
\hline
\texttt{correct\_dedup} & Duplikat binarny lub semantyczny wskazujący
oczekiwany obraz. \\
\hline
\texttt{wrong\_dedup} & Duplikat binarny lub semantyczny wskazujący
inny obraz. \\
\hline
\texttt{correct\_substitute} & Podmiana oczekiwanego obrazu. \\
\hline
\texttt{wrong\_substitute} & Podmiana innego obrazu. \\
\hline
\texttt{miss} & Obraz przyjęty jako nowy (\texttt{inserted}). \\
\hline
\texttt{failed} & Błąd przetwarzania. \\
\hline
\end{tabular}
\end{table}

Notatnik rysuje mapy cieplne liczności wyników w grupach i etykiet
oraz rozkład odległości dla każdej etykiety. Zbiorcze metryki
zapisuje do pliku \texttt{metrics.yaml}. W repozytorium notatnik
nie zawiera zapisanych wyników, dlatego w tym podrozdziale opisano
wyłącznie metodę.

\subsubsection{Ograniczenia procedury}

Analiza kodu notatników wskazuje warunki, które należy uwzględnić
przy interpretacji wyników:

\begin{itemize}
    \item \textbf{Brak zapytań rozpraszających.} Wybierane są tylko
    zapytania z przyporządkowaniem. Procedura nie mierzy więc
    fałszywych alarmów dla obrazów bez odpowiednika, które według
    autorów DISC2021 obniżają \(\mu\mathrm{AP}\)~\cite{douze2021disc}.
    \item \textbf{Obecność obrazu referencyjnego.} Zapytania i obrazy
    referencyjne losowane są niezależnie. Kod nie sprawdza, czy obraz
    oczekiwany dla zapytania znalazł się wśród 200 przesłanych obrazów
    referencyjnych. Gdy go brak, przyjęcie zapytania jako nowego jest
    poprawne, a mimo to otrzymuje etykietę \texttt{miss}. Przed
    obliczeniem czułości należy ograniczyć zapytania do tych, których
    obraz referencyjny jest w próbie.
    \item \textbf{Duplikaty binarne.} Brama ustawia
    \texttt{existing\_image\_key} tylko w potoku podobieństwa.
    Zdarzenie \texttt{hash\_duplicate} nie wskazuje więc istniejącego
    obrazu, a notatnik porównuje pustą nazwę z oczekiwaną i przypisuje
    etykietę \texttt{wrong\_dedup}. Poprawnie wykryte kopie dokładne
    są w tej grupie liczone jako błędy.
    \item \textbf{Kolejność przetwarzania.} Lista jest mieszana,
    a w partii działa 10 równoległych zadań. Kopia może zostać
    przetworzona przed oryginałem. Wtedy kopia otrzymuje wynik
    \texttt{inserted}, a oryginał staje się duplikatem. Etykiety
    przypisywane są jednak tylko kopiom, więc taki przypadek liczony
    jest jako \texttt{miss}.
    \item \textbf{Liczność próby.} Przy 200 zapytaniach brak błędów
    danego rodzaju ogranicza ich odsetek z góry tylko do około
    1,5\%~\cite{hanley1983rule}.
\end{itemize}

Strojenie progu wymaga więc dwóch pomiarów: rozkładu odległości par
z oznaczeniem \(P\) i \(N\) na danych zbliżonych do docelowych oraz
oceny całego potoku z uwzględnieniem obrazów rozpraszających.
Pierwszy pomiar wskazuje zakres progu z równania~\eqref{eq:tpr-fpr},
drugi sprawdza wpływ przybliżonego wyszukiwania, kolejności przesyłania
i podmiany reprezentanta na wynik końcowy.
